\subsection{Описание метода}

\subsubsection{Интерполяционный многочлен в форме Лагранжа}

Рассмотрим произвольный многочлен степени \( N \):
\[
  L_N(t)=\sum_{n=0}^N a_nt^n
\]

Чтобы он был интерполяционным, необходимо условие:
\[
  L_N(t_n)=f_n,\quad n=\overline{0,N}
\]

Оно равносильно решению СЛАУ:
\[
  \begin{cases}
    a_0+a_1t_0+\ldots+a_Nt_0^N=f_0\\
    a_0+a_1t_1+\ldots+a_Nt_1^N=f_1\\
    \dots\\
    a_0+a_1t_N+\ldots+a_Nt_N^N=f_N\\
  \end{cases}
\]

Определим совместность СЛАУ:
\[
  \Delta = \begin{vmatrix}
  1 & t_0 & t_0^2 & \dots & t_0^N \\
  1 & t_1 & t_1^2 & \dots & t_1^N \\
  \vdots & \vdots & \vdots & \ddots & \vdots \\
  1 & t_N & t_N^2 & \dots & t_N^N
  \end{vmatrix} = \prod_{i \neq j}^{N} (t_i - t_j), \quad 0 \le j < i \le N
\]

Матрица невырождена, если узлы интерполяции попарно различны. Допустим, это так. Значит, решение есть и оно единственно.

По методу Крамера найдём искомый многочлен:
\[
  L_N(t)=\sum_{n=0}^N\dfrac{\Delta_n}{\Delta}\cdot t^n
\]

Хотя интерполирующая функция найдена, для представления в форме Лагранжа разложим \( \Delta_n \) по элементам \( n \)-го столбца:
\[
  a_n=\dfrac{\sum_{i=0}^Nf(t_i)\Delta_{ni}}{\Delta}
\]

Перегруппируем слагаемые, чтобы собрать члены с одинаковыми \( f(t_n) \):
\[
  L_N(t)=\sum_{n=0}^Nf(t_n)\cdot\Phi_n(t)
\]

Здесь в качестве \( \Phi_n(t) \) обозначены некоторые линейные комбинации \( \{t^i\}^N_{i=0} \). Можно заметить, что они обладают свойством:
\[
  \Phi_i(t_j)=\delta_i^j,
\]

где \( \delta_i^j \) --- символ Кронекера.

Из этого следует, что многочлены имеют вид:
\[
  \Phi_i(t) = A(t - t_0)(t - t_1) \dots (t - t_{i-1})(t - t_{i+1}) \dots (t - x_n)
\]

Зная, что \( \Phi_i(t_i)=1 \), имеем:
\[
  \begin{gathered}
    1 = A(t_i - t_0)(t_i - t_1) \dots (t_i - t_{i-1})(t_i - t_{i+1}) \dots (t_i - t_n)\implies\\
    \implies\Phi_i(t) = \frac{(t - t_0)(t - t_1) \dots (t - t_{i-1})(t - t_{i+1}) \dots (t - t_n)}{(t_i - t_0)(t_i - t_1) \dots (t_i - t_{i-1})(t_i - t_{i+1}) \dots (t_i - t_n)}
  \end{gathered}
\]

Таким образом, мы вывели формулу для \textit{интерполяционного многочлена в форме Лагранжа}:
\[
  L_N(t)=\sum_{n=0}^N f(t_n)\cdot\prod_{\substack{j=0\\ n\neq j}}^N\dfrac{t-t_j}{t_n-t_j}
\]

\subsubsection{Многочлен Чебышёва}

Функция вида
\[
  T_n(t)=\cos(n\arccos t),\quad t\in[-1,1],\ n\in\mathbb{N}_0
\]

называется \textbf{многочленом Чебышёва первого рода}.

Доказательство факта, что приведённая функция является многочленом степени \( n \), выполняется методом математической индукции.

Найдём нули функции \( T_n(t) \):
\[
  \begin{gathered}
    \cos(n\arccos t)=0\implies n\arccos t_k=\dfrac{\pi}{2}+\pi k,\quad k\in\mathbb{Z}\implies\\
    \implies\arccos t_k=\dfrac{2k+1}{2n}\pi\implies t_k=\cos\left(\dfrac{2k+1}{2n}\pi\right),\quad\dfrac{2k+1}{2n}\pi\in[0,\pi]
  \end{gathered}
\]

Поскольку это многочлен степени \( n \), он имеет не более \( n \) различных вещественных корней \textit{(основная теорема алгебры)}. Из \( \cos\colon\mathbb{R}\to\mathbb{R} \) следует, что все \( n \) корней лежат на отрезке \( [-1,1] \):
\[
  t_k=\cos\left(\dfrac{2k-1}{2n}\pi\right),\quad k=\overline{1,n}
\]

Масштабируем корни на произвольный отрезок \( [a,b] \):
\[
  t_k=\dfrac{a+b}{2}+\dfrac{b-a}{2}\cos\left(\dfrac{2k-1}{2n}\pi\right)
\]

Таким образом, мы получили набор интерполяционных узлов \( \{t_n\}_{n=0}^N\subset[a,b] \) и можем приступить к составлению численного метода.